% Chapter 7

\chapter{Conclusions}\label{Chapter7}

%----------------------------------------------------------------------------------------
%	SECTION 1
%----------------------------------------------------------------------------------------

\section{Summary of Key Findings}
The purpose of this study was to find out if data-driven clustering of neuropsychological testing could show a cognitive pattern in acquired brain injury (ABI) that would be reliable, understandable and useful to clinicians. Four hypotheses were tested using a pipeline that consisted of UMAP dimensionality reduction, HDBSCAN density-based clustering, and multiple imputation methods for missing clinical data.



Overall, the major finding of this research was that the cognitive structure found through the clustering process represents a cognitive severity gradient, not a distinct number of phenotypes. The cohort represents a single dimension of overall cognitive severity. The first principal component accounts for approximately 56\% of the total variability among the six cognitive domains, with almost identical positive weights. The pipeline takes this continuum and divides it into three tiers: Above Average (38.6\%, \(n = 6{,}725\)), Near Normal (36.4\%, \(n = 6{,}328\)), and Globally Impaired (25.0\%, \(n = 4{,}353\)). The level determines which tier a participant falls into. However, there is no evidence that shape affects this outcome.



Hypothesis 1 is supported with an important caveat. There are identifiable clusters that are relatively stable. However, these are best viewed as severity bands. Each of the ten imputation methods meets the pre-specified silhouette and noise thresholds, and the two extreme tiers remained fairly consistent under Hennig per-cluster matching (Core Jaccard 0.63 and 0.62, respectively; §\ref{sec:h1_results}).



Hypothesis 2 is also supported. Using a common pipeline, the average Adjusted Rand Index across all forty-five pairwise method comparisons is .710 (95\% bootstrap CI [.704, .716]), greater than the .500 threshold. Following Holm-Bonferroni adjustment, forty-four of forty-five pairwise combinations reject \(H_0: \mathrm{ARI} \leq .5\) at \(\alpha = .05\). Only KNN-SoftImpute is borderline. Consensus clustering achieves a mean confidence of .929, with 86.2 percent of assessments receiving high-confidence assignment. The baseline comparison with listwise deletion verifies that imputation is analytically required: removing missing values from 46.9 percent of the assessments alters the number of clusters and lowers the silhouette value to .257.



Hypothesis 3 is supported as well. Clinical Unit is significantly correlated with cognitive severity tier ($\chi^2 = 911.3$, $p \approx 4 \times 10^{-156}$); however, the magnitude of the correlation is small (Cramér's \(V = .162\)). The Cochran-Mantel-Haenszel statistic demonstrates that clinical unit is not solely determined by diagnosis group \parencite{cochran1954some}. Thus, cognitive severity appears to contribute to clinical unit designation independently of diagnostic category.



Lastly, Hypothesis 4 is supported. Domain-level aggregation generates cluster separation comparable to variable-level features (silhouette .481 vs. .473), and domain-level aggregation also decreases multivariate collinearity (mean VIF 1.88 vs. 2.77) and increases numerical stability (11.8 vs. 61.7). These results provide empirical support for conceptualizing cognitive performance at the domain level.

\section{Methodological Contribution}
There are two methodological contributions made within this dissertation. First, an imputation-robustness framework is presented. The combination of ARI/NMI comparison and consensus confidence measures provides a systematic means to evaluate whether differences in clustering result arise due to differing missing-value algorithms. Due to its domain-agnostic nature, the framework may be employed with other types of partially observed clinical or biomedical data, including genomic data or survey responses.



Second, a cautionary note is provided regarding feature engineering. As demonstrated in this dissertation, unsupervised cognitive phenotyping is very susceptible to how features are constructed. For example, inclusion of an unintended demographic variable or failure to include any cognitive test information in the assessment record can create an artifact representing what appears to be a phenotypic grouping that in reality arises from poor construction of features. The pipeline developed in this dissertation reduces this potential source of error by excluding non-cognitive variables, ensuring that at least one cognitive variable exists for each assessment record, and performing Winsorisation, direction-alignment, and standardization on each variable prior to aggregation. Clinical labeling should occur only after these processing steps have taken place and evaluation for sensitivity.

\section{Clinical Recommendations}
Based upon these findings, I recommend the following:



\begin{itemize}

    \item Use a global severity index to supplement a diagnosis. The global severity index provides a reproducible summary description of overall cognitive function at time of assessment, regardless of the type of assessment (including imputed values for uncompleted tests). Rehabilitation teams can use both the diagnosis and a global severity index to inform their determination of appropriate treatment intensity for each patient.

\end{itemize}



\begin{itemize}

    \item Label uncertain cases instead of discarding them. Consensus clustering provides high confidence labels for 86.2 percent of assessments but about 2{,}400 assessments lie close to a tier boundary and thus will vary depending on choice of algorithm. Decision-support systems should report the confidence with which labels were generated instead of artificially confining label certainty uniformly.

\end{itemize}



\begin{itemize}

    \item Choose model-based imputation for practical use. Although MICE, EM, MissForest, and autoencoder imputation produce similarly robust results (pairwise ARI values as large as .88) compared to each other, KNN tends to behave differently from the others and Mean Imputation loses variance relative to the others. Therefore, MICE seems like a reasonable choice for practical application purposes whereas MissForest may prove useful for nonparametric modeling applications where nonlinear relationships are expected.

\end{itemize}



\begin{itemize}

    \item Develop treatment objectives at the domain level. Clustering at the domain level produces separations among clusters similar to variable-level aggregation while simultaneously reducing multivariate collinearity (mean VIF 1.88 vs. 2.77) and increasing numerical conditioning (11.8 vs. 61.7). This lends further support to the long-established neuropsychological practice of evaluating batteries according to domain of cognitive functioning \parencite{lezak2012neuropsychological}.

\end{itemize}

\section{Final Remarks}
Using a dimensional reduction technique called UMAP and clustering technique called HDBSCAN on reduced data with imputation algorithms that handle missing values, this study demonstrates that unsupervised clustering can uncover a continuous cognitive severity gradient in ABI that is robust to different imputation strategies and clinical groups and interpretable at the domain level. Furthermore, this research shows that there are no separate phenotypes in this cohort but rather a single dimension of overall cognitive severity.



Therefore, this dissertation provides empirical support for developing rehabilitation planning guidelines using severity levels for patients who have suffered an acquired brain injury. Future studies will be needed to validate the usefulness of using cognitive severity levels for rehabilitation planning purposes:



\begin{itemize}

    \item Does cognitive severity levels predict future outcomes better than diagnosis and current measures?

\end{itemize}



\begin{itemize}

    \item Would using a severity system improve resource allocation in controlled environments or is it just a formalization of what was discovered in CMH analyses?

\end{itemize}



\begin{itemize}

    \item Are there additional subgroups present in independent cohorts across multiple centers that were not evident in this sample or do they emerge due to varying characteristics of the populations sampled and/or assessment times?

\end{itemize}
